\addcontentsline{toc}{chapter}{Abstract}

\chapter*{Abstract}

\vspace*{-8mm}

Physics-Informed Neural Networks with sinusoidal activation functions
(SIREN-PINNs) are applied to simulate optical speckle patterns by solving the
2D Helmholtz equation with a complex electric field and a random-phase boundary
condition. The proposed SIREN architecture ($5 \times 128$, $\omega_0 = 1.0$)
achieves relative $L^2$ errors of $0.006\%$ (1D) and $0.171\%$ (2D) against
analytical solutions, several orders of magnitude below the $5\%$ thesis
threshold.
Speckle patterns are generated by imposing a random phase
$\phi(x) \sim \mathcal{U}(0, 2\pi)$ on the illuminated surface and solving the
Helmholtz equation freely in the interior domain.
Statistical validation confirms fully developed speckle: contrast
$C = 1.0253$ and negative-exponential intensity distribution consistent with
Goodman's theory.
A two-phase Adam~+~L-BFGS optimization strategy and Latin Hypercube Sampling
(LHS) are key to convergence.
Reproducibility was verified with three independent random seeds, yielding a
mean $L^2$ error of $0.155 \pm 0.020\%$.

\bigskip

\noindent\textbf{Keywords:} Physics-Informed Neural Networks, SIREN,
Helmholtz equation, optical speckle, sinusoidal activation, Latin Hypercube
Sampling.
